\subsection{E5: Scalability and Cost}
\label{sec:e5}

\paragraph{Setup.}
We evaluate end-to-end reconstruction cost and scalability on a fixed sample of
30 schemata (10~RSchema, 10~SchemaBench, 10~Enterprise-145; seed~42), stratified
by size and constraint complexity.
Each gold DDL is verbalized by GPT-5.4-mini or Claude Sonnet~4.6 under the exact
E1 basic/detailed prompts, then reconstructed by a fixed GPT-4o-mini constructor
(E2 protocol).
We report micro-averaged F1, mean nGED, sequential wall-clock latency, and
provider-usage USD ($N{=}10$ matched schemata per cell; 120/120 pipelines completed).

\paragraph{Results.}
Table~\ref{tab:e5_cost} summarizes the sample.
Under basic prompting, referential F1 is high for both describers
(RSchema: GPT 0.909 / Claude 0.971; SchemaBench: 0.968 / 0.995;
Enterprise-145: 0.960 / 0.960), while Domain F1 remains weaker on SchemaBench
(${\approx}0.53$--$0.61$).
Claude costs ${\approx}1.9$--$3.0\times$ GPT-5.4-mini on matched cells
(mean round-trip \$0.038 vs.\ \$0.018).
Detailed prompting raises cost in every cell
($\Delta\$ {\approx}0.004$--$0.021$) but does \emph{not} improve referential F1
on this sample ($\Delta$Ref from $-0.052$ to $-0.233$).
Larger schemata drive higher latency and cost (Enterprise Large median latency
often ${>}250$\,s); the single RSchema Large schema ($N{=}1$) is not treated as
a stable size estimate.
On the 10 overlapping DeepSeek-basic IDs, GPT-5.4-mini basic exceeds DeepSeek
referential F1 (0.909 vs.\ 0.667 on RSchema; 0.968 vs.\ 0.629 on SchemaBench),
with DeepSeek monetary cost unavailable.
